\subsection{权与特殊权}

命题 3. \(W_{a}\) 的特殊点（第五章第 3 节第 10 号定义 1）正是 \(R^{\prime}\) 的权。

设 \(x_0 \in E\) 且 \(\alpha \in R\)。平行于 Ker \(\alpha\) 且经过 \(x_0\) 的超平面 L 的方程为 \(\langle \alpha, x \rangle = \langle \alpha, x_0 \rangle\)。要使它等于某个 \(L_{\beta,k}\)，一方面必须 \(\alpha\) 与 \(\beta\) 成比例，因 R 是约化的，故 \(\beta = \pm \alpha\)；另一方面必须 \(\langle \alpha, x_0 \rangle\) 为整数。立即得到，\(x_0\) 是 \(W_a\) 的特殊点，当且仅当 \(\langle \alpha, x_0 \rangle \in Z\) 对所有 \(\alpha \in R\) 成立，换言之，当且仅当 \(x_0 \in P(R^{\prime})\)（第 1 节，第 9 号）。

推论。(i) 若 \(\overline{\omega} \in \mathrm{P}(\mathbb{R}^{\prime})\)，则存在一个小室 C，使得 \(\overline{\omega}\) 是 \(\overline{\mathbb{C}}\) 的一个极点。

\begin{enumerate}
\def\labelenumi{(\roman{enumi})}
\setcounter{enumi}{1}
\tightlist
\item
  若 \(\mathbf{C}\) 是一个小室，则 \(\overline{\mathbf{C}} \cap \mathbf{Q}(\mathbb{R}^{\prime})\) 退化为一个点，且该点是 \(\overline{\mathbf{C}}\) 的一个极点。
\end{enumerate}

这由命题 3、第五章第 3 节第 10 号命题 11 的推论以及第五章第 3 节第 10 号命题 12 可得。

命题 4. 设 \(C'\) 是 \(R'\) 的一个室。

\begin{enumerate}
\def\labelenumi{(\roman{enumi})}
\item
  存在唯一的小室 C 包含于 \(\mathbf{C}'\)，使得 \(0\in \overline{\mathbf{C}}\)
\item
  \(w(\overline{\mathbf{C}})\)（\(w\in \mathrm{W}\)）的并，是 E 中 0 的一个邻域。
\item
  \(\mathbf{C}'\) 的每一面壁都是 \(\mathbf{C}\) 的一面壁。
\end{enumerate}

这由第五章第 3 节第 10 号命题 11 可得。

现在假设 \(R\) 不可约。设 \((\alpha_i)_{i\in I}\) 是 \(R\) 的一组基（第 1 节，第 5 号，定义 2），并设 \((\overline{\omega}_i)_{i\in I}\) 是对偶基。\(\overline{\omega}_i\) 是 \(R^{\prime}\) 的基本权，属于室 \(D'\)；此室是 \(R^{\prime}\) 中对应于基 \((\alpha_i)\) 的室。令

\[
\tilde {\alpha} = \sum_ {i \in \mathrm{I}} n _ {i} \alpha_ {i}
\]

为 \(\mathbf{R}\) 的最高根（第 1 节，第 8 号），并令 \(\mathbf{J}\) 为由 \(i \in \mathbf{I}\) 且满足 \(n_i = 1\) 的 i 所成的集合。

命题 5. 设 \(\mathbf{C}\) 是包含于 \(\mathbf{C}'\) 且其闭包包含 0 的小室（命题 4）。

\begin{enumerate}
\def\labelenumi{(\roman{enumi})}
\tightlist
\item
  C 是由 \(x \in E\) 组成的集合，其中 \(\langle\alpha_{i}, x\rangle > 0\)（对所有 \(i \in I\)），且有 \(\langle\tilde{\alpha}, x\rangle < 1\)。
\item
  集合 \(\overline{C} \cap P(R^{\prime})\) 由 0 以及 \(\overline{\omega}_{i}\)（\(i \in J\)）组成。
\end{enumerate}

令 D 为由 \(x \in E\) 组成且满足 \(\langle\tilde{\alpha}, x\rangle < 1\) 的集合，并置 \(C_{1} = C' \cap D\)。由于 \(0 \in \overline{C}\)，有 \(C \subseteq D\)，从而 \(C \subseteq C_{1}\)。我们将证明，对于所有 \(\alpha \in R\) 和 \(k \in Z\)，集合 C 与 \(C_{1}\) 位于超平面 \(L_{\alpha,k}\) 的同一侧。这将证明 \(C_{1} \subseteq C\)，从而建立断言 (i)。若 k = 0，则整个室 \(C'\) 位于 \(L_{\alpha,0}\) 的一侧，在此情形断言成立。若 \(k \neq 0\)，将 \(\alpha\) 替换为 \(-\alpha\) 后可设 k \textgreater{} 0。于是，在 C 上有 \(\langle\alpha, x\rangle < k\)，这是因为 \(0 \in \overline{C}\)。另一方面，\(\tilde{\alpha} - \alpha\) 在 \(C'\) 上为正（第 1 节，第 8 号，命题 25）。因此，对 \(y \in C_{1}\)，有 \(\langle\alpha, y\rangle \leqslant \langle\tilde{\alpha}, y\rangle < 1 \leqslant k\)。所以 C 与 \(C_{1}\) 位于 \(L_{\alpha,k}\) 的同一侧。

现在设 \(\overline{\omega} \in P(R^{\prime})\)。则 \(\overline{\omega} = \sum_{i} p_i \overline{\omega}_i\)，其中 \(p_i \in \mathbf{Z}\)（第 1 节，第 10 号），并且 \(\overline{\omega} \in \overline{C'}\) 当且仅当整数 \(p_i\) 为正。若 \(\overline{\omega} \in \overline{C'}\)，则 \(\overline{\omega} \in \overline{C}\) 当且仅当 \(\langle \tilde{\alpha}, \overline{\omega} \rangle = \sum_{i} n_i p_i\) 满足 \(\leqslant 1\)，故得 (ii)。

推论。小室 C 是一个开单纯形，其顶点为 0 以及 \(\overline{\omega}_i / n_i\)（\(i \in I\)）。

这由 (i) 得出。

